\section{Scope and reproducibility}\label{sec:scope}

The theorem concerns a fixed compact Kähler background, fixed locally
free factors, two nonzero layers, and fixed multiplicities. Its
uniformity includes approaches to faces through the simple stable
region; it does not include the nonsimple endpoint itself. Higher
extensions have composable arrows and nonlinear integrability
relations, so the exact harmonic presentation and parity argument
used here do not apply without further work. Singular graded sheaves
and actual Yang--Mills flow estimates require additional analysis.

\subsection{Complex endomorphisms and adjoint symmetry}
Let $L_Z$ be the finite operator associated with
$\mathcal Q_{s,Z}$ and the rank mass. In orthonormal arrow
coordinates, direct expansion gives
\begin{equation}\label{eq:adjointsymmetry}
r_v\bigl((L_Z(X^*))_v-(L_ZX)_v^*\bigr)
=[X_v^*,\mu_s(Z)_v].
\end{equation}
At central balance the right side vanishes. Hence the Hermitian
real fixed space of adjoint complexifies to the full operator,
with the same eigenvalue multiplicities. This agrees with the
centralizer identity of~\cite[\S5.2, Proposition~5.3]{hkkp2023}.
Thus the six-mode example is a failure of scalar-vertex compression,
not a failure of correctly complexifying the Hermitian operator at
central balance.

Off balance a Hermitian restriction can lose even the correct
kernel. Take one target and one source, each of dimension two,
with arrows $I_2$ and
$J=\left(\begin{smallmatrix}0&1\\0&0\end{smallmatrix}\right)$.
The identity arrow forces the two centralizer blocks to agree;
commutation with $J$ then gives $X=Y=aI+bJ$. On the full
trace-free complex space the kernel is the line $X=Y=bJ$.
Its Hermitian part is zero. This representation is not centrally
balanced: its moments are $\operatorname{diag}(2,1)$ and
$\operatorname{diag}(-1,-2)$. It illustrates why the intermediate
gauge comparison in Lemma~\ref{lem:gaugespectrum} is performed
on the complex space. Likewise a single identity arrow has
connected support but all simultaneous matrix endomorphisms in
its centralizer. Connectedness alone never replaces the
scalar-centralizer condition.

\subsection{Finite calculations}
The accompanying scripts use the explicit matrices printed in
Section~\ref{sec:example}. Exact rational computations verify the
covariances, full centralizer dimension, positive trace-free Gram
form, characteristic polynomials, and the off-balance kernel
comparison above. A separate implementation forms the raw complex
commutator matrix in a Pauli basis and verifies the same six-mode
polynomial and the coefficient $25/20$. These calculations check
finite algebra; the proofs of HE existence and relative elliptic
estimates are those in Sections~\ref{sec:curvature}--\ref{sec:spectrum}.

For the plotted family we solve the two real target balance equations.
This is valid for these real matrices: uniqueness of the positive
Hermitian minimizer and invariance under complex conjugation imply
that it is real symmetric. Small eigenvalues are evaluated as the
squares of singular values of the commutator map restricted to a
mass-orthonormal trace-free basis. This avoids forming a rounded
Gram matrix before computing its weakest eigenvalue.
The numerical coefficient at $\delta=10^{-10}$ is approximately
$1.24999999996$, consistent with the proved limit $5/4$.

The source distribution contains the calculation scripts, generated
data, figure source, dependency versions, and commands for rebuilding
the figure and this PDF. No partial differential equation is solved
numerically in these computations. The numerical output is an
illustration and reproducibility check of the finite model, not an
additional hypothesis in the theorems.
